\section{PROYECTO}

El proyecto consiste en un \textbf{tacómetro digital con control de velocidad} para un
motor de corriente continua con reductora, implementado sobre la placa de desarrollo
OM13085 (LPCXpresso1769) con el microcontrolador LPC1769 de NXP.

Sobre el eje del motor se monta un disco de 60 ranuras que interrumpe el haz de un
sensor óptico de ranura. Contando los flancos que produce el sensor durante una ventana
de tiempo fija, el sistema calcula la velocidad de giro en RPM. A partir de esa medición
el motor puede operar en dos modos, que el usuario elige con una llave selectora:

\begin{itemize}
  \item \textbf{Modo manual:} la posición de un potenciómetro (perilla) fija directamente
  el ciclo de trabajo de la señal PWM que alimenta al motor, de 0 a 100\,\%.
  \item \textbf{Modo automático:} el usuario envía desde una computadora, por puerto serie,
  una consigna de velocidad entre 0 y 185 RPM. Un regulador proporcional-integral (PI)
  ajusta el PWM para que la velocidad medida siga a la consigna.
\end{itemize}

En paralelo, el sistema mide la temperatura ambiente con un sensor LM35. Si la
temperatura supera los 45\,\textdegree C se activa una alarma visual (LED rojo) y acústica
(buzzer), que se apaga recién cuando la temperatura baja a 43\,\textdegree C o menos.

Cada 1s el equipo envía a la computadora una trama de telemetría con la velocidad
medida, la temperatura y el estado de operación, de modo que el usuario puede supervisar
el funcionamiento desde una terminal serie.

\subsection{OBJETIVOS}

\begin{itemize}
  \item Medir la velocidad del eje con un sensor óptico y un disco de 60 ranuras,
  usando un temporizador del microcontrolador como contador de eventos.
  \item Accionar el motor mediante PWM, con mando manual por potenciómetro o regulación
  automática a partir de una consigna recibida por UART.
  \item Adquirir las señales analógicas (potenciómetro y LM35) con el ADC del
  microcontrolador, transfiriendo los resultados a memoria mediante DMA.
  \item Señalizar con prioridad una condición de sobretemperatura.
  \item Informar periódicamente el estado del sistema a una computadora.
  \item Organizar el firmware en capas (aplicación, servicios y dispositivos), de manera
  que la lógica de la aplicación no dependa del hardware concreto.
\end{itemize}

\subsection{REQUISITOS}

\begin{longtable}{|p{1.2cm}|p{12.4cm}|}
\hline
\textbf{ID} & \textbf{Requisito} \\ \hline
\endhead
R01 & Ejecutar el firmware sobre el LPC1769 de la placa OM13085. \\ \hline
R02 & Medir la velocidad en RPM con un disco de 60 ranuras y un sensor óptico. \\ \hline
R03 & En modo manual, el potenciómetro determina el ciclo de trabajo del PWM del motor. \\ \hline
R04 & En modo automático, regular la velocidad según una consigna recibida por UART. \\ \hline
R05 & Las fuentes de mando son excluyentes: en cada modo actúa una sola. \\ \hline
R06 & Al entrar en modo automático el motor queda con PWM cero hasta recibir una consigna
nueva y válida; no se reutilizan consignas anteriores al cambio de modo. \\ \hline
R07 & Si la temperatura supera el umbral, encender el LED rojo y el buzzer. \\ \hline
R08 & Adquirir potenciómetro y temperatura con ADC y DMA. \\ \hline
R09 & Por UART, recibir consignas y transmitir velocidad, temperatura y estado. \\ \hline
R10 & Firmware en tres capas; la aplicación accede al hardware sólo mediante servicios. \\ \hline
R14 & Aceptar consignas de 0 a 185 RPM inclusive (velocidad máxima del motor con PWM al
100\,\%). \\ \hline
\end{longtable}

\subsection{ALCANCE}

El trabajo comprende el diseño del hardware de interconexión de los sensores y
actuadores con la placa OM13085, y el desarrollo completo del firmware en lenguaje C. El
firmware usa la biblioteca CMSIS y los drivers LPC17xx provistos por la cátedra sin
modificarlos, y se compila tanto con MCUXpresso IDE como con CMake y la cadena GNU Arm.

La interfaz de usuario en la computadora es una terminal serie estándar; no se desarrolla
software específico para la PC. El sistema no detiene el motor por sobretemperatura: la
alarma es sólo de señalización.
